\chapter{The plug march}
\label{ch:plug}

Every nozzle in this document so far has been a bell: a solid tube
that wraps the gas on all sides. This chapter builds the other kind.
A plug nozzle keeps solid metal only \emph{under} the flow and lets
the outside of the jet touch the atmosphere directly --- one change
that buys the nozzle the ability to adapt itself to its
surroundings, and costs the solver a boundary it has never had to
handle. The step adds that boundary, assembles a complete plug
march around it, and certifies the assembly against a flow whose
answer is known exactly on paper.

\section{The question}

The deepest results in the RDE nozzle theory are about plugs
(the corpus's plug theorem, its T4)~\cite{theory} --- in particular
about what a plug does when the chamber pressure swings over a
cycle, which is Chapter~\ref{ch:plugcycle}. None of them can
be tested until the engine can march a plug at all. So: what does
the solver need that it does not yet have, and can the new assembly
be checked against an exact answer rather than only against itself?

\section{How it works, in full}

\subsection{The topology, and the one new piece of physics}

A plug (or \emph{aerospike}) nozzle is a bell turned inside out. In
a bell the gas is enclosed by wall all the way to the exit. In a
plug the wall lies only below the flow --- a central body, the
\emph{spike} --- and above the flow there is nothing but the
surrounding atmosphere. The outer surface of the exhaust is
therefore a \emph{free jet edge}: a boundary of the gas that no
hardware holds in place.

That is the whole appeal of the shape. Because the outer boundary
is free, the jet widens or narrows on its own until its pressure
matches the atmosphere. A bell has one fixed exit area, so it is
correctly expanded at exactly one ambient pressure and wrong at
every other; a plug re-shapes its own jet as the ambient pressure
falls with altitude. That is the classical
\emph{altitude-compensation} argument for plugs
(Chapter~\ref{ch:nozzle}) --- and it is why plugs matter here for a
second reason: an RDE swings its chamber pressure every cycle,
which the nozzle feels the same way it feels a change of altitude.

The free edge is the one new piece of physics the solver must
learn. On a wall, the \emph{position} is known and the pressure
comes out as a result; on a free edge it is exactly the reverse.
The pressure is prescribed --- $p = p_a$, the ambient pressure,
because the atmosphere pushes back and nothing else does --- and
the \emph{position} is what gets solved. In the cell language of
Chapter~\ref{ch:moc}, the free-edge cell is fed by the previous
edge point (the edge is a \emph{streamline}: a line the flow runs
along and no gas crosses) and by the $C^{+}$ characteristic rising
from the interior point below it in its own column. The prescribed
pressure fixes the speed modulus $q_{p_a}$ through the gas tables,
so the two remaining unknowns are the new position and the turning
angle $\theta_4$, with
$(u,v) = q_{p_a}(\cos\theta_4, \sin\theta_4)$
(Figure~\ref{fig:freejet}). The assembly then mirrors the bell
march upside down: spike wall below (the inverse-wall cell of
Chapter~\ref{ch:wallinput}, mirrored to the $C^{-}$ family that a
bottom wall receives), free edge above, no axis~\cite{raospike}.

\begin{figure}[htbp]
\centering
\includegraphics[width=0.95\textwidth]{fig_freejet.pdf}
\caption{The free-boundary cell, drawn in the physical plane:
$x$ axial [m] against $y$ radial [m]. (a)~The new edge point pt4 is
built from two known states --- the previous edge point pt3, whose
flow direction sets the direction of the edge segment because the
edge is a streamline, and the interior neighbour pt1 in the same
column, which sends the $C^{+}$ characteristic. What to look at is
which quantity is held and which floats: the pressure is
\emph{imposed}, $p = p_a$ exactly, which fixes the speed modulus
$q_{p_a}$ through the tables, while the position $(x_4,y_4)$ and
the turning angle $\theta_4$ are solved. That inversion --- the
opposite of a wall cell, where position is given and pressure comes
out --- is what makes a plug a plug. (b)~The same cell with the
sign trap it survived on the way to certification: writing the
boundary condition as $v = +\sqrt{q_{p_a}^2-u^2}$ admits only
outward turns (Appendix~\ref{app:signtrap}).}
\label{fig:freejet}
\end{figure}

\subsection{Why the lip fan is constructed, not marched}

At the \emph{cowl lip} --- the sharp outer edge of the channel that
feeds the plug --- the gas leaves the internal duct and turns
around a corner into the open. A supersonic gas turning away from
itself expands through a \emph{Prandtl--Meyer fan}: a continuous
sheaf of straight rays, each carrying one intermediate state,
across which the flow accelerates, its pressure falls, and its
direction rotates smoothly from the incoming angle to the outgoing
one. When the corner is sharp, every ray starts from the same
geometric point and the fan is \emph{centered} on it. That point is
\emph{singular}: all the intermediate states coexist there at once,
so the flow has no single value at the lip and its gradients are
infinite (Figure~\ref{fig:fansing}).

This violates exactly the assumptions an ordinary cell rests on. A
cell whose segment spans the fan closes its compatibility relation
with a midpoint average of two states that differ by the whole turn
--- an order-one error, not a second-order one --- and near the lip
all the $C^{-}$ characteristics meet at one point, so the position
crossing degenerates. Measured: a certification residual of
$8.9\times10^{9}$ when an early build tried exactly this.

The cure costs nothing, because a centered fan is self-similar: the
state on each ray depends only on the ray angle and is given
exactly by the Prandtl--Meyer function. So the corner is never
marched, it is \emph{constructed} --- seed the fan as data, then
resume ordinary cells between rays, where the flow is smooth again.
Brick 1 does precisely this at its sharp throat corner (the initial
expansion fan); the oracle below goes one step further and places
the start line downstream of the lip with the exact fan field
written on it, so the singular point lives in the data and never
inside a cell.

\begin{figure}[htbp]
\centering
\includegraphics[width=\textwidth]{fig_fan_singularity.pdf}
\caption{Why the lip fan cannot be marched through. All panels are
the physical plane, $x$ [m] against $y$ [m]. (a)~The lip as a
singular point: a uniform stream at $M_1$ and a uniform stream at
$M_2$ meet across a centered fan whose every ray emanates from one
geometric point, so at the lip every state between $M_1$ and $M_2$
coexists and the gradient is infinite. (b)~An ordinary cell asked
to step across it: its segment spans the fan, so the midpoint
closure averages two states separated by the whole turn, and the
$C^{-}$ crossings degenerate as they converge on the lip ---
measured, a certification residual of $8.9\times10^{9}$. (c)~The
cure: each ray is seeded with its exact Prandtl--Meyer state and
ordinary cells resume between rays; the oracle puts the start line
downstream, carrying the exact field. What this proves is that the
singularity is handled by construction, not by refinement --- no
mesh is fine enough to march through a point where the solution is
multivalued.}
\label{fig:fansing}
\end{figure}

\subsection{The march: columns, growth at the edge, consumption at
the wall}

Each column is swept \emph{bottom-up}: first the wall point at the
prescribed station, then the interior points, each fed by the
$C^{+}$ from the point just solved below it and the $C^{-}$ from
its partner on the previous column, and finally the free-edge point
on top, fed from its own column (why it must be its own column is a
lesson the oracle taught: Appendix~\ref{app:edgeroot}).

The wall cell finds its foot --- the point on the previous column
where the characteristic arriving at the wall started --- by a
\emph{multi-column search}: it scans the whole previous column,
stepping further columns back if needed, for the segment whose data
brackets its arriving $C^{-}$ (the bracketing function is positive
at the previous wall point and decreases with row height, so the
first sign change marks the foot segment). This lets the march run
at fine prescribed stations with short compatibility chords; the
earlier scheme, which confined the foot to the first segment of the
adjacent column, had to stretch its columns and paid for it in
accuracy (Appendix~\ref{app:footbias}).

Row bookkeeping closes the topology, one characteristic at a time,
at both boundaries. At the \emph{edge}, each column's free-edge
point becomes a new top row: the widening jet gains one crossing
per column (the mirrored $j_2{+}{=}1$ of the bell march; what
happens without it is Appendix~\ref{app:pluglessons}, lesson 2).
At the \emph{wall}, rows are instead \emph{consumed}: the incoming
$C^{-}$ family terminates on a solid wall, so any previous-column
rows lying below the new wall point's foot have already met the
wall between stations and must not be prolonged --- the new
column's interiors are partnered with the previous rows \emph{from
the foot up}, re-indexed. Prolonging them was the last defect the
oracle caught: the terminated characteristics kept marching as
ghost rows below the spike --- for such a row the only crossing
Newton can still find is the backward one, a valid root of the
equations on the wrong branch --- and the mass check exposed the
ghost limb at $+9\%$. With consumption, mass closes to 0.026\%.

\subsection{The oracle}

The march is certified against a flow whose \emph{entire field} is
known in closed form, so the marched answer can be compared with
the truth rather than with itself. The flow is a planar corner jet:
uniform $M_1 = 1.5$, a centered fan at the lip, uniform $M_2 = 2.2$
turned by $\Delta\nu = 23.9^{\circ}$, with the ambient set to
$p_a = p(M_2)$. The prescribed spike is the exact streamline of
that field, obtained by fine-step RK4 integration --- an
independent reference, not another method-of-characteristics run
--- and the start line carries the exact field on a vertical cut.

The marched solution must then reproduce four things, each within a
tolerance derived from the scheme's own convergence order rather
than chosen by hand (Chapter~\ref{ch:method}): the edge's
asymptotic angle and speed (P-1), the wall pressure distribution
station by station (P-2), mass conservation through the final
column (P-3), and a \emph{momentum balance} (P-4). The last one is
worth spelling out. Draw a closed \emph{control surface} around the
gas --- an imaginary boundary whose faces are the start line, the
final column, the spike surface and the free edge --- and require
that the momentum flowing out equals the momentum flowing in plus
the push the pressures exert on the faces, every pressure measured
relative to ambient. On the free edge the pressure \emph{is}
ambient, so measured that way the free face pushes exactly zero:
the plug's defining property, here tested rather than assumed.
Alongside these, a deliberately corrupted ambient (R-1) must leave
the bands --- a check that the checks can fail. Every cell is
Newton-certified.

\subsection{The axisymmetric stage}

For the axisymmetric plug on the real-gas tables no closed form
exists, so that stage certifies by \emph{internal} closures ---
conservation laws the solution must satisfy whatever the answer is.
The start-line data is a tables-consistent corner fan
($\mathrm{d}\theta = \sqrt{M^2-1}\,\mathrm{d}q/q$ integrated on the
tables' own $M(q)$), the prescribed spike is a streamline of that
data, and the checks are: Newton certification (A-1); mass
conservation with the $2\pi y$ flux weight of an axisymmetric ring
(A-2); the same ambient-gauge momentum balance (A-3); a
\emph{source-alive} control --- the same data marched with the
axisymmetric source term switched off ($\delta = 0$, planar) must
differ by more than the axisymmetric pair's own tolerance band,
which guards against an axisymmetric mode that is silently running
planar (A-4); and a corrupted-spike rejector (A-5).

\section{What was built}

\carrier{validation/a1\_freejet\_unit.py}: the free-jet cell alone,
against unit known answers. \carrier{validation/a1\_plug\_march.py}:
the assembly --- bottom-up columns at prescribed stations, the
mirrored inverse-wall cell, the multi-column foot search, row
growth at the edge and row consumption at the wall --- plus the
conservation instruments (\code{col\_fluxes},
\code{wall\_push\_poly}), the closed-form oracle with its rejector
(\code{STAGE=oracle}), and the axisymmetric internal-closure stage
(\code{STAGE=axi}).

\section{Results of record}

\textbf{The free-jet cell: 6/6 PASS.} A uniform jet already at
ambient pressure is an exact fixed point; a turned jet is preserved
exactly in planar flow --- in axisymmetric flow the source term
legitimately changes a turned uniform state, a physics point the
first test tripped over; a compressing edge turns inward; and the
Prandtl--Meyer corner known answer lands at machine precision ---
edge speed vs.\ closed form $1.9\times10^{-9}$, edge angle equal to
the exact turn angle to $1.9\times10^{-16}$.

\textbf{The assembly vs the oracle: 6/6 PASS} (66~s, stations
41/81, start rows 21/41). Every cell Newton-certified --- worst
residual 0.019 (coarse, 1253 cells) and 0.017 (fine, 4870 cells) of
the allowed tolerance. P-1: the marched edge's tail angle is
0.41809 vs.\ the exact $\Delta\nu = 0.41777$ --- $|d| =
3.2\times10^{-4}$ against a band of $4.5\times10^{-3}$, with the
edge speed equal to the closed form to all printed digits. P-2:
98\% of wall stations inside the derived band, maximum deviation
$9.8\times10^{-3}\,p_a$. P-3: mass through the final column 8869.98
vs.\ start-line 8867.65 --- $|d| = 2.33$ (0.026\%) against a band
of 26. P-4: momentum closure $6.9\times10^{3}$ on
$F_{\rm in} = 2.0\times10^{7}$ ($3.4\times10^{-4}$ relative; band
$7.2\times10^{4}$; coarse/fine ratio 3.6, confirming the scheme's
second order), with the free edge contributing zero by
construction. R-1: a 1\% corruption of the ambient moves the edge
angle $20\times$ outside the P-1 band --- rejected.
Figure~\ref{fig:plug} shows the certified march; the road here ---
four structural lessons, each measured --- is
Appendix~\ref{app:pluglessons}.

\textbf{The axisymmetric stage: 5/5 PASS} (39~s, stations 33/65,
start rows 17/33, NASA real-gas tables, tables-fan turn
0.4395~rad). A-1: all cells Newton-certified, worst residuals
0.017/0.022. A-2: mass with the $2\pi y$ weight, $|d| = 505$ on
$1.085\times10^{5}$ (0.47\%) against a band of 986. A-3: momentum
closure $-2.4\times10^{6}$ on $F_{\rm in} = 2.4\times10^{8}$
($1.0\times10^{-2}$ relative, inside its band). A-4: the
axisymmetric source is alive --- the wall speed at the last station
differs from the planar run by 156~m/s, $9\times$ the band. A-5: a
corrupted spike slope breaks the momentum closure at $5\times$ the
band --- rejected.

\begin{figure}[htbp]
\centering
\includegraphics[width=\textwidth]{fig_plug.pdf}
\caption{The certified march against the closed-form oracle.
(a)~The physical plane, $x$ [m] against $y$ [m]: the exact world
--- spike (red), exact jet edge (blue dashed), lip fan rays (grey),
start line (dotted) --- overlaid with the marched characteristic
mesh (light dots) and the marched free edge (orange). Look at how
the orange edge tracks the blue line: told only the pressure on
that surface, the solver found where the surface is. (b)~Wall
pressure along the spike, $x$ [m] against $p_w$ [bar]: marched
values against the closed form, with the derived tolerance band ---
98\% of stations inside it, worst deviation
$9.8\times10^{-3}\,p_a$. The two panels together are the
certificate: the free surface and the solid surface are both right.
Compare with the state before the fixes,
Figure~\ref{fig:plugbefore}.}
\label{fig:plug}
\end{figure}
